\documentclass[11pt,a4paper]{article}

\begin{document}

\tableofcontents

\section{M6-7}

\subsection{a}

$\overrightarrow{F_g}=-G\frac{m_em_p}{d^2}\overrightarrow{u_r}$\\
$\overrightarrow{F_e}=\frac{1}{4\pi\epsilon_0}\frac{q_eq_p}{d^2}\overrightarrow{u_r}$\\
Or $Gm_em_p\approx10^{-69}$ et $\frac{1}{4\pi\epsilon_0}e^2\approx10^{-24}$\\
Les interactions gravitationelles sont donc complètement négligeables.

\subsection{b}

On utilise le théorème du moment cinétique.\\
$\frac{d\overrightarrow{L_O}}{dt}=M_O(\overrightarrow{F_e})=\overrightarrow{OP}\wedge\overrightarrow{F_e}=\overrightarrow{0}$\\
Donc $\overrightarrow{L_O}$ est constant, donc $\overrightarrow{OP}$ et $\overrightarrow{v}$ sont toujours contenus dans le même plan.

\subsection{c}

D'après le principe fondamental de la dynamique, on a :\\
$m_ea=F_e=\frac{1}{4\pi\epsilon_0}\frac{q_eq_p}{mr^2}$\\
$2r\frac{d^2\theta}{dt^2}\overrightarrow{u_\theta}-r\frac{d\theta}{dt}^2=\frac{1}{4\pi\epsilon_0}\frac{-e^2}{m_er^2}\overrightarrow{u_r}$\\


\end{document}